%!TEX root = main.tex
\chapter{Driefasen}
\label{Driefasen}

Bij het gebruik van drie fase zijn er 2 verschillende schakelingen mogelijk. Y-Y of $\Delta -\Delta$.

\section{Y-Y configuratie}

Dit is de schakeling voor de Y-Y of ster-ster configuratie:

\begin{center}
    \begin{circuitikz}
        \draw (0,0) to[isourcesin, l_=$V_{bn}$] (2,1.75)
        to [short, o-o, i_=$I_{bB}$] (3,1.75)
        to [european resistor, l_=$Z_y$] (5,0)
        to [european resistor, l_=$Z_y$] (5,-2)
        to [short, o-o, i<_=$I_{cC}$] (0,-2)
        to [isourcesin, l_=$V_{cn}$] (0,0);
        \draw (0,0) to[isourcesin, l_=$V_{an}$] (-2,1.75)
        to [short] (-2,3)%7,1.75
        to [short, o-o, i>_=$I_{aA}$] (7,3)
        to [short] (7,1.75) %-2,3
        to [european resistor, l^=$Z_y$] (5,0); 
        \draw[dotted] (5,0) to[short, i_=$I_{Nn}$] (0,0);
    \end{circuitikz}
\end{center}

Voor de spanningen geldt het volgende:

$V_{an}=V_y\angle 0^\circ$

$V_{bn}=V_y\angle -120^\circ$

$V_{cn}=V_y\angle 120^\circ$

\vspace{5mm}

De gestippelde line is de line to neutral stroom.

Voor de fase spanningen geldt:

$I_{aA}=\frac{V_y\angle 0}{Z_y\angle \theta ^\circ}=I_L\angle -\theta ^\circ$

$I_{bB}=\frac{V_y\angle -120}{Z_y\angle \theta ^\circ}=I_L\angle -120^\circ -\theta ^\circ$

$I_{cC}=\frac{V_y\angle 120}{Z_y\angle \theta ^\circ}=I_L\angle 120^\circ -\theta ^\circ$

Hierbij geldt ook, $I_{aA}+I_{bB}+I_{cC}=0$, en dus $I_{nN}=0$ 

\vspace{5mm}

Voor het berekenen van het vermogen worde de volgende formule gebruikt:

$p(t)=3V_y\cdot I_L\cdot \cos \theta$

\vspace{5mm}

Lijnspanning: de spanning tussen 2 lijnen.

$V_{ab}=V_{AB}=$lijnspanning

$V_{ab}=V_{an}-V_{bn}$

$V_{an}=V_{ab} \cdot \frac{1}{\sqrt{3}}\angle -30^\circ $

$V_{ab}=V_{an} \cdot \sqrt{3} \angle 30^\circ$

\section{$\Delta -\Delta $ configuratie}

Dit is de schakeling voor de $\Delta -\Delta $ / delta - delta of driehoek - driehoek, configuratie.

\begin{center}
    \begin{circuitikz}
        \draw (0,0) to[isourcesin] (3,0)
        to [isourcesin] (1.5,-2.25)
        to [isourcesin] (0,0);
        \draw (5,0) to[european resistor, l^=$Z_\Delta $] (8,0)
        to [european resistor, l^=$Z_\Delta $] (6.5,-2.25)
        to [european resistor, l^=$Z_\Delta $] (5,0);
        \draw (3,0) to[short, o-o, i<_=$I_{bB}$] (5,0);
        \draw (1.5,-2.25) to[short, o-o, i>_=$I_{cC}$] (6.5,-2.25);
        \draw (0,0) to[short] (0,1)
        to[short, o-o, i>_=$I_{aA}$] (8,1)
        to[short] (8,0);
    \end{circuitikz}
\end{center}

In deze configuratie is geen line tot neutral te vinden. Voor de spanningen geldt hey volgende:

$V_{ab}=$ lijnspanning

$V_{ab}=V_L\angle 30^\circ$

$V_{bc}=V_L\angle -90^\circ$

$V_{ca}=V_L\angle 150^\circ$

\vspace{5mm}

De fasestroom: stroom door de fase (belasting).

$V_{ab}=V_{AB}$

$I_{AB}=\frac{V_{AB}}{Z_\Delta}=\frac{V_{ab}}{Z_\Delta}=\frac{V_L \angle 30^\circ}{Z\angle \theta^\circ}=I_A \angle 30^\circ - \theta ^\circ$

$I_{CA}=\frac{V_{CA}}{Z_\Delta}=\frac{V_{ab}}{Z_\Delta}=\frac{V_L \angle 150^\circ}{Z\angle \theta^\circ}=I_A \angle 150^\circ - \theta ^\circ$

$I_{aA}=I_{AB}-I_{CA}$

$I_{aA}=\sqrt{3}\cdot I_{AB} \angle -30^\circ$

\vspace{5mm}

Als we een belasting toevoegen op de lijn aA kan dit niet berekend worden. Hier voor zou de driehoek naar de ster configuratie
worden om gezet. Hierbij geldt:

$V_{an}=V_{ab}\cdot \frac{1}{\sqrt{3}}\angle -30^\circ$

$Z_\Delta =3Z_y$

\vspace{5mm}

Bij het berekenen van het vermogen geldt: 

$U_{load}\cdot I_{load} \cdot \cos \varphi$

$U_{load}\cdot I_{load} \cdot \sin \varphi$

$\varphi$ kan je direct uit de hoek van de impendantie halen.

\vspace{5mm}

De vermogens van bijde schakelingen, dus Y-Y en $\Delta - \Delta$, zijn altijde het zelfde.

